رتن ناتھ سرشار کی ہزار داستان

دوسرا حصہ
حجام اور قزاق

بچوں کیلئے


تسہیل و تلخیص
اقبال کاردار



فیروز سنز پرائیویٹ لمیٹیڈ
لاھور۔ راولپنڈی۔ کراچی





راز کھلتا ہے
اگلی شب ملکہ شہرزاد نے اَدھوری کہانی یوں شروع کی:
"اے شاہِ جہاں پناہ! حسن بدرالدّین کو کچھ دیر کے لیے یہیں چھوڑتی ہوں اور دلھن اور کبڑے کا حال سناتی ہوں۔"
صبح ہوئی تو دلھن یہ دیکھ کر حیران رہ گئ کہ اس کے دُلھا حسن بدرالدین کا کہیں پتا نہیں۔ اتنے میں اس کا باپ وزیر شمس الدّین آیا کہ چل کر دیکھوں بیٹی کس حال میں ہے۔ دروازے پر پہنچ کر اپنی لڑکی کا نام پکارا تو وہ مارے خوشی کے جامے میں پھولی نہ سمائی۔ اٹھکیلیاں کرتی ہوئی نکل آئی۔"ابّا جان بندگی!" کہ کر زمیں کو بوسہ دیا۔
"جانِ بابا! تم اس قدر مسرور کیوں ہو؟"
 اس پر وزیر زادی بولی "بابا جان! خاتونِ جنت کی قسم کھا کر کہتی ہوں کہ میری شادی بظاہر تو اس کبڑے سے ہوئی تھی لیکن اصل میں میرا دولھا کوئی اور تھا۔ کبڑے سے شادی صرف دکھاوا تھی۔ آپ بھی شکر کیجیے اللہ نے ہماری عزت رکھ لی۔"
وزیر اور بھی حیران ہوا ، بولا "جانِ بابا! میں یہ کیا سن رہا ہوں! مجھے اپنے کانوں پر یقین نہیں آتا۔"
وزیرزادی نے ابھی جواب نہیں دیا تھا کہ ساتھ والے کمرے سے کسی کے کراہنے کی آواز سنائی دی۔ وزیر فوراً اس کمرے میں گیا تو کیا دیکھتا ہے کہ وہی کبڑا ایک جگہ اوندھا پڑا ہواہے۔ ٹانگیں اُوپر اور سر نیچے۔ یہ عجیب تماشا دیکھ کر وزیر شمسُ الدّین دنگ رہ گیا۔ اُدھر کبڑا یہ سمجھا کہ  جِن پھر کمرے میں آگیا ہے۔ اس نے کراہنا بند کر دیا اور یوں دم سادھ لی گویا زبان پر قفل پڑگیا ہو۔
وزیر نے پوچھا "کیوں  بے!یہ کیا بیہودگی ہے؟ تیرا دماغ تو نہیں چل گیا ہے!" مگر کبڑےنے کوئی جواب نہ دیا۔ کیوںکہ اُسے معلوم تھا کہ اگر اُس نے زبان ذرا بھی ہلائی تو جِن اسے جیتا نہ چھوڑے گا۔ وزیر بہت جھلّایا اور گرج کر بولا "اگر تو نے جواب نہ دیا اور یوں ہی دم سادھے پڑا رہا تو ابھی تیری گردن اُڑا دوں گا۔"
اِس پر کبڑےنے خوف سے کانپتے ہُوئے جواب دیا۔ "اے جِنوں کے سردار! خُدا کی قسم کھا کر کہتا ہُوں کہ جِس وقت سے تو نے میری ٹانگیں اُوپر اور سر نیچے کیا ہے، جب سے اِسی حالت میں پڑا  ہوں۔ ذر ا بھی نہیں ہلا ہوں۔ میں تیرے آگے ہاتھ جوڑتا ہوں کہ مجھ غریب مِسکین پر رحم کر اور  اِس سزا سے نجات دے۔ میں وعدہ کرتا ہوں کہ پھر کبھی اِدھر کا رُخ نہ کروں گا۔"
وزیر کی حیرت اور بڑھی۔ یا خُدا یہ کیا اِسرار ہے؟ یہ کمبخت مُجھے جِنوں کا سردار سمجھ بیٹھا ہے۔ پھر جھلّا کر پوچھا:
" کمبخت ! یہ کیا بکتا ہے؟ میں جِنوں کا سردار نہیں وزیر شمسُ الدّین ہوں۔ دولھن کا باپ۔"
لیکن کبڑے نے جواب دیا۔" اے شخص! جا اپنا کام کر اور راہ لے! مجھے یوںہی پڑا رہنے دےکہ میری جان اُس جِن کے ہاتھ میں ہےجو مجھے اِس حالت میں چھوڑ کر گیا ہے۔وہ مجھے کسی جِن کی لڑکی یا بھینس سے بیاہ دے گا اور میرا خانہ خراب کرے گا۔خُدا اُس کو روزِ بد دکھائے اور میرے حال کو پہنچے۔"
وزیر بولا "میں تجھے حکم دیتا ہوں کہ اُٹھ اور یہاں سے دفع ہو جا۔ پھر کبھی منحوس صُورت نہ دکھانا۔ نہیں تو کھال اُدھیڑ دوں گا۔ "
کبڑا ٹَس سے مَس نہ ہُوا اور بولا "آپ کا حکم سر آنکھوں پر۔ لیکن بندہ جِن کی اِجازت کے بغیر ہرگز نہیں جائے گا۔ میرا اُس کا عہد ہے۔ اگر میں نے وعدہ کی خلاف ورزی کی تو جِن کو کیا منہ دکھاوں گا؟ وہ مجھے حُکم دے گیا ہے کہ سورج نکلنے کے بعد یہاں سے جانا۔ لہٰذا بندہ مجبور ہے۔"
تب وزیر نے پوچھا ۔"اچھا یہ بتا ، تو کس کے ساتھ یہاں آیا تھا؟ تیرے ساتھ دوسرا  نوجوان کون تھا؟ جھوٹ مت بولنا، نہیں تو ہڈی پسلی ایک کر دوں گا۔"
کبڑے نے جواب دیا ۔"میرے ساتھ ایک خوبصورت نوجوان یہاں آیا تھا۔ شکل و صورت اور لباس سے کوئی رئیسذادہ لگتا تھا۔"
پِھر کبڑے نے اَپنی ساری سرگزشت اوّل تا آخر سُنا دی ۔"جب میں اِس کمرے میں پہنچا تو اچانک پانی میں سے دھواں اُٹھا اور ہوتے ہوتے ایک بھینس میں تبدیل ہو گیا۔ اِس نے مجھے وہ ڈرایا وہ ڈرایا کہ میرے ہوش اُڑ گئے۔ جاتے وقت میرا سر نیچے اور ٹانگیں اوپر کر گیاکہ  صبح تک یوں ہی پڑے رہنا، نہیں تو کچا چبا جاؤں گا۔خُدا اِس دولھن کو غارت کرے کہ جس کے سبب میری گت بنی ہے۔"
کبڑے کی یہ بات سُن کر وزیر کو بہت غصہ آیا، کبڑے کے دو تین لاتیں جمائیں اور حکم دیا ۔"ابھی دفع ہو جا یہاں سے۔ گھوم کر دیکھا بھی تو آنکھیں نکلوا دوں گا۔"
کبڑا روتا ہوا یہاں سے چلا اور سلطان کے دربار میں جا کر اپنی آپ بیتی کہہ سنائی۔ بادشاہ نے کہا ۔"چل دور ہو یہاں سے۔ تیرا مقدر ہی خراب ہے ورنہ آج تو وزیر شمسُ الدّین کا داماد ہوتا اور یوں تیری ٹانگیں اُوپر اور سر نیچے نہ ہوتا۔"
اب وزیر شمس الدّین کا حا ل سُنیے۔
اِس نے جو کبڑے بدبخت کی کہانی سنی تو بہت حیران ہوا کہ یہ کیا راز ہے؟ پھر اپنی بیٹی سے پوچھا "جانِ بابا! اِس کبڑے نے جو کہانی سنائی ہے اِس میں کیا راز ہے؟ یہ جِن اور بھینس کا کیا قصہ ہے؟ اور وہ نوجوان کون تھا جو کبڑے کے ساتھ یہاں آیا اور صبح ہوتے ہی غائب ہو گیا؟"
وزیر کی بیٹی نے جواب دیا "بابا جان! کبڑے کا حال میں کیا جانوں؟ ہاں اِس نوجوان کے بارے میں کہہ سکتی ہوں کہ وہ بہت شریف اور کوئی امیرزادہ لگتا تھا۔ نجانے وہ رات کے کس پہر کو یہاں سے چلا گیا۔ یہ پگڑی اِس کی ہے جو آپ کے سامنے رکھی ہے۔"
یہ سُن کر وزیر نے پگڑی اُٹھائی اور بغور دیکھنے لگا۔ "یا خُدا!  یہ تو کسی وزیر کی پگڑی ہے اور مُوصل شہر کی بنی ہوئی ہے۔" اِس کے بعد ایک تعویذ پگڑی کی ٹوپی کی تہہ میں سِلا ہوا پایا۔ شوق چرایا کہ کھل کر دیکھوں تو سہی اِس میں کیا چیز ہے! کھولا تو اِس تعویذ میں کاغذ نظر آیا جس میں یہودی کا عہد نامہ تحریر تھا۔ حسن بدرالدّین ابن نورالدّین باشندہِ قاہرہ(مصر) کا نام پڑھا۔ اِس کے بعد  ایک اور کاغذ پڑھا تو زور کی چیخ ماری اور غش کھا کر زمین پر گر پڑا۔
لڑکی نے باپ کی جو یہ حالت دیکھی تو گبھرا گئ۔ صراحی سے پانی کا گلاس بھرا اور باپ کے منہ پر پانی کے چھینٹے دیے۔
چند منٹ بعد وزیر نے ہوش میں آ کر آنکھیں کھول دیں۔ پھر بلند آواز سے کہا "واہ باری تعالیٰ تیرا کوئی شریک نہیں۔ تو کریم و کارساز ہے۔ قادرِمطلق اور بندہ نواز ہے۔" پھر اَپنی بیٹی کو کاغذ پر لکھی ہوئی تحریر پڑھ کر سنائی اور بولا:
"بیٹی! کچھ سمجھی؟ کس  کے ساتھ تمھارا نکاح ہوا ہے۔"
وہ بولی "بابا جان! میں اِس کے نام سے آگاہ نہیں ہوں۔ مجھے تو بس اِتنا ہی پتا ہے کہ وہ اِس کبڑے کے ساتھ یہاں آیا تھا۔"
تب وزیر نے بتایا۔
"بیٹی! تیرا  میاں، میرے چھوٹے بھائی  نورالدّین کا بیٹا ہے۔ اِس کا نام حسن بدرالدّین ہے۔ واقعی تو بڑی خوش قسمت ہے بیٹی! خُدا نے بڑا کرم کیا۔ ﷲ اکبر اﷲ اکبر ! اِس کی کریمی کے قربان! سمجھ  میں نہیں آتا ، یہ لڑکا یہاں کیونکر  آیا؟ کِس فرشتے نے اِسے یہاں پہنچایا مجھے کانوں کان خبر نہیں۔"
پھر باقی تحریر پڑھی تو نورالدّین کی شادی کی تاریخ درج پائی۔ اور پھر حیران ہوا کہ یہ وہی تاریخ تھی جب اُس کی اپنی شادی ہوئی تھی۔ اِس کے بعد حسن بدرالدّین کی پیدائش کی تاریخ لکھی تھی۔ پڑھتے ہی غنچہِ دل کھِل اُٹھا۔ اِسی تاریخ کو اِس کی بیٹی پیدا ہوئی تھی۔
یہ سب کاغذات لے کر وزیر شمسُ الدّین بادشاہ کے دربار میں حاضر ہوا اور تخت کو بوسہ دے کر کہا "جہاں پناہ! یہ عجیب اتفاق ہے کہ میرا اور نورالدّین کا نکاح ایک ہی تاریخ کو ہوا اور ایک ہی تاریخ کو میرے ہاں لڑکی اور اُس کے ہاں لڑکا پیدا ہوا اور اَب حُسنِ اِتفاق سے اُسی کے لڑکے  حسن بدرالدّین سے میری بیٹی کا نکاح ہو گیا۔ حالاں کہ ہمیں کانوں کان خبر نہیں۔"
بادشاہ یہ سن کر بڑا حیران ہوا۔ حکم دیا ۔"یہ عجیب و غریب واقعہ لکھ لیا جائے کہ اِس طرح کی داستان اِس سے پہلے سننے میں نہیں آئی"۔ اب وزیر اپنے بھتیجے اور داماد  حسن بدرالدّین کی آمد کا منتظر رہا اور جب عرصے تک کوئی خبر نہ پائی تو بہت فکرمند ہوا۔ کچھ سوچ کر دل میں کہا کہ اَب میں وہ کام کروں گا جو کسی سے نہیں ہوا تھا اور نہ آج تک کسی کو سوجھا تھا۔





حسن بدرالدین کا بیٹا
 اب سُنیےکہ کچھ عرصے بعد حسن بدرالدّین کی بیوی یعنی وزیر  شمسُ الدّین کی بیٹی کے ہاں چاند سا لڑ کا پیدا ہوا۔ وزیر شمسُ الدّین نے نواسے کی پیدائش پر بہت خوشی منائی۔ غریبوں کو عمدہ کھانا کھلایا۔ مُحتاجوں میں کپڑے بانٹے اور دوست و احباب کی دعوت کی۔ کافی سوچ بچار کے بعد اپنے نواسے کا نام عجیب رکھا۔
عجیب کی عمر بارہ برس کی ہو گئ مگر اِس کے باپ  حسن بدرالدّین کا کچھ پتا نہ چلا کہ کہاں اور کِس حال میں ہے۔ شمسُ الدّین  کو بڑی تشویش ہوئی ۔پھر دل میں سوچا کیوں نہ بادشاہ سے اجازت لے کر حسن کی تلاش میں جاؤں۔ بیٹی سے مشورہ کیا تو  اِس نے بھی یہی  کہا کہ ہمیں اُس کو تلاش کرنا چاہیے۔ تب وزیر نے دربار میں جا کر بادشاہ سے کہا۔ بادشاہ نے بخوشی اِجازت دے دی۔ اِشارے کی دیر  تھی، سب سامان مہیا  ہو گیا۔ وزیر نے اپنی بیٹی اور نوسے عجیب کو ساتھ لیا اور سفر پر چل نکلا۔
تیسرے روز  یہ لوگ شہر دمشق میں اور کچھ دن تک یہیں ٹھہرنے کا ارادہ کیا۔ ایک کشادہ میدان میں خیمے لگائے۔ وزیر نے اپنے خادموں سے کہا ۔"دو دن تک ہمارا یہاں قیام رہے گا۔ جس کسی کو جو دیکھنا ہو دیکھ آئے، جو لینا ھو لے آئے۔" 
خادموں نے جو یہ سنا تو شہر  کی سیر کو گئے۔ کوئی اِس غرض سے گیا کہ کچھ خرید لائے، کوئی اِس لئے گیا کہ کچھ فروخت کر کے نئی چیزیں خرید لائے، کوئی حمام کو گیا اور کوئی دمشق کی مشہور مسجد کی زیارت کو چلا۔ عجیب بھی ایک خواجہ سِرا کو ہمرا لے کر شہر کی سیر کو گیا۔
شہر کے لوگوں نے جو عجیب کو دیکھا تو عَش عَش کر اُٹھے۔ ایسا خوبرُو لڑکا اُنہوں نے کم ہی دیکھا تھا۔ جِدھر  سے وہ نکل جاتا، ایک زمانہ جمع ہو جاتا۔ ایک پر ایک گِرا پڑتا۔ 
حُسنِ اتفاق سے ایک غلام حسن بدرُالدّین کی دُکان کے سامنے ٹھر گیا۔ یہ وہی دُکان تھی، جس میں قاضی اور گواہوں کے سامنے اُس ڈاکو باورچی نے حسن بدرُالدّین کو اپنا بیٹا بنایا تھا۔ باورچی کو مرے ھوئے اب عرصہ ھو چکا تھا اور حسن بدرُالدّین اُس کی تمام جائداد کا اکیلا وارِث تھا۔ ایک غلام کے وہاں ٹھرنے سے دوسرے غُلام بھی وہیں ٹھر گئے۔ اچانک حسن بدرُالدّین کی نظر  عجیب پر پڑی۔ دیکھ کر عَش عَش کر اُٹھا۔ اُس کا جی چاہا کہ ابھی دکان سے اُتر کر اِس لڑکے کو اَپنے گلے سے لگا لے۔ اُس کا دل بےقابو ہو کر عجیب کی طرف کھِنچا جا تا تھا اور وہ حیران تھا کہ یا الہٰی اِس لڑکے کی محبت میں میرا خون کیوں اِسقدر جوش کھانے لگا ہے دوسری حیران کُن بات یہ تھی کہ حسن بدرُالدّین اور عجیب کی صورت ایک دوسرے سے بہت مِلتی جُلتی تھی۔ 
 بدرُالدّین نے اُسی وقت اَنار کا شربت تیار کیا اور  لڑکے کی طرف دیکھ کر بولا۔ "اے آقا زادہ! اگر مضائقہ نہ ہو  تو یہ شربت نوش فرمائے۔"
           عجیب نے جو یہ اِلفاظ سُنے تو اُس کا دل بھر آیا۔ خواجہ سِرا کی طرف دیکھ کر کہا۔"اِس شخص کی باتوں نے میرے دل پر جادو کا کام کیا ہےاور میرے دل میں بھی اِس  کے لیے محبت پیدا ہوگئ ہے۔ خُدا جانے یہ کون ہے! مگر مُجھے اِس سے ہمدردی پیدا ہوگئ ہے۔ میرا خیال ہے اس کا کوئی لڑکا گم ہوگیا ہے۔چلو اِس کے دل کو تسلی دیں اورجو کچھ کِھلائے، خوش ہو کر کھا آئیں۔ شاید اِس طرح سے خُدا میرے بچھڑے ہُوۓ باپ کو مجھ سے مِلا دے اور  ہمارے سوئے ہوۓ بھاگ جھاگ اُٹھیں۔"
        مگر خواجہ سِرا نے منع کیا اور کہا ۔"حضور! اﷲ جانتا ہے يہ بات نامناسب ہو گی۔ آپ تو لڑکے بن کر چھوٹ جائیں گے مگر غلام کو سزا بُھگتنا پڑے گی"۔ مگر عجیب کہاں ماننے والا تھا۔ خواجہ سِرا کو ساتھ لے کر دُکان میں آ گیا۔ بدرُالدّین نے انار  کا میٹھا اور خوشبودار  شربت پیش کیا۔ خواجہ سِرا اور عجیب نے مل کر پیا۔
جب وہ دونوں شربت پی کر فارغ ہوئے تو حسن بدرُالدّین نے شکریہ ادا کیا کہ آپ نے غریب خانے پر قدم رنجہ فرمایا۔ عِزّت بَخشی مرتبہ بڑھایا۔ اَب آپ کچھ تناول فرمائیں۔ اِس پر عجیب نے کہا ۔"باورچی تم بھی ھمارے ساتھ کھانے میں شریک ہو جاؤ۔ شائد اِس بہانے سے بچھڑے ہوؤں کو  خُدا ایک بار پھر ملا دے"۔
حسن بدرُالدّین نے دریافت کیا ۔"کیا حضور کا بھی کوئی عزیز جدا ہو گیا ہے؟" 


